\documentclass[10pt,a4paper]{amsart}
\usepackage[margin=19mm]{geometry}
\usepackage{amsmath,amssymb,mathtools}
\usepackage[scr=rsfs,cal=euler]{mathalfa}
\usepackage{xfrac}
\usepackage[unicode,hidelinks,bookmarks=false]{hyperref}
\newcommand{\ie}{i.e.\ }
\newcommand{\cf}{cf.\ }
\newcommand{\resp}{respectively\ }
\newcommand{\Subsection}{\subsection}
\providecommand{\nobreakdash}{\nobreak\hbox{-}}
\setlength{\emergencystretch}{2em}
\multlinegap0pt
\usepackage{amssymb}
\usepackage{mathtools}
\usepackage{xcolor}
\newtheorem{pro}{Proposition}
\newtheorem{thm}{Theorem}
\newcommand{\A}{\widehat{A}}
\def\Alex{\operatorname{Alex}}
\newcommand{\Alpha}{\widehat{\alpha}}
\def\Aut{\operatorname{Aut}}
\newcommand{\C}{\mathbb{C}}
\newcommand{\Id}{\operatorname{id}}
\renewcommand{\phi}{\varphi}
\title{Fourier Analysis and Idempotents in Quandle Algebras}
\author{Mohamed Elhamdadi, Luc Ta, Bryce Virgin}
\date{Source: arXiv:2609.03799v1 (CC BY 4.0)}
\begin{document}
\maketitle

\noindent\emph{Source and licence.} Fourier Analysis and Idempotents in Quandle Algebras. Version arXiv:2609.03799v1 (CC BY 4.0), distributed by arXiv under the Creative Commons Attribution 4.0 International licence, http://creativecommons.org/licenses/by/4.0/, as recorded in the arXiv licence metadata for this exact version and shown on the abstract page https://arxiv.org/abs/2609.03799v1. Full licence notice: \texttt{licenses/arxiv-2609.03799-CC-BY-4.0.txt}.\par\smallskip

\noindent\textit{Editorial scope.} Counted unit: the article's thm thm (source0.tex lines 300--306). The article's complete original proof of it is delivered with it (source0.tex lines 308--320, one \texttt{proof} environment of 13 lines). No mathematics is written, repaired or re-derived: the statement and the proof are the article's own lines, in the article's own notation, and no retained line is substituted or re-flowed.  Retained uncounted as the source-local context the proof needs: lines 267--274.  Nothing was omitted from any retained span, so the package carries no omission record.  No retained line contains a citation, so the wrapper carries no bibliography and no bibliography entry.  No retained line contains a figure environment, an image-inclusion command or drawing code, so no image is delivered and the package declares no asset dependency.  Editorial formatting only, in addition: an AMS \texttt{amsart} preamble that restates the article's own declarations and macros used by the retained lines, namely the environments \texttt{pro}, \texttt{thm} on separate counters, together with the standard packages those lines need.  The retained environments print in this document as Pro 1, Theorem 1 in order of appearance, whereas the article prints the same items with the numbering of its own full document.  The complete native source of the article as distributed in the arXiv e-print for this exact version is bundled unchanged as \texttt{sources/209315/source0.tex}.
\begin{pro}\label{pro:l1}
    For each $\phi\in\Aut(A)$, the algebra $L^1(A,\mu,\C)$ has a well-defined product
    \begin{align*}
        (\alpha \beta)(x) &:= \Delta_A\phi^{-1}\int_{ A}(\alpha \circ S_y^{-1})(x) \beta(y)d\mu(y)\\
        &= \Delta_A\phi^{-1}\int_{ A}\alpha (\phi^{-1}(x)+(\Id-\phi^{-1})(y) ) \beta(y)d\mu(y).
    \end{align*}
    If moreover $\Alex(A,\phi)$ is latin, then in fact \[ (\alpha \beta)(x) = \Delta_A(\Id-\phi)^{-1} \int_{ A}\alpha (y')   \beta(L_{y'}^{-1}(x))d\mu(y'). \]
\end{pro}

\begin{thm} 
    Let $\Alex(A,\phi)$ be a $\sigma$-compact, locally compact latin Alexander quandle, and let $\alpha,\beta\in L^1(A,\mu,\mathbb{C}) $. Then
    \[
    \widehat{\alpha \beta} (\chi)=\Alpha (\chi\circ\phi)\hat{\beta} (\chi\circ(\Id-\phi))
    \]
    for all $\chi\in\A$.
\end{thm}

\begin{proof}
    Fix $\chi\in\A$. Using the definition of the Fourier transform and Proposition \ref{pro:l1}, we compute
    \begin{align*}
        \widehat{\alpha\beta}(\chi)&=\int_{A}(\alpha\beta)(x)\overline{\chi(x)} d\mu(x) \\
        &=\Delta_A(\Id-\phi)^{-1} \int_{A}\int_{ A}\alpha (y)   \beta(L_{y}^{-1}(x))d\mu(y)\overline{\chi(x)} d\mu(x) \\
        &=\Delta_A(\Id-\phi)^{-1} \int_{A}\int_{ A}\alpha (y)   \beta(L_{y}^{-1}(x))d\mu(y)\overline{\chi(x)} d\mu(\phi(y)+(\Id-\phi)(L_y^{-1}(x)))\\
        &=  \int_{A}\int_{ A}\alpha (y)   \beta(L_{y}^{-1}(x))d\mu(y)\overline{\chi(\phi(y))}\overline{\chi((\Id-\phi)(L_y^{-1}(x)))} d\mu(  L_y^{-1}(x))\\
        &=  \int_{A}\int_{ A}\alpha (y)   \beta(z)d\mu(y)\overline{\chi(\phi(y))}\overline{\chi((\Id-\phi)(z))} d\mu(  z)\\
        &=\left(\int_{  A}\alpha (y)\overline {\chi(\phi(y))} d\mu(y)\right) \left(\int_{  A}\beta(z)\overline {\chi((\Id-\phi)(z))} d\mu(z)\right)\\
        &=\Alpha(\chi\circ\phi)\hat{\beta}(\chi\circ(\Id-\phi)).
    \end{align*}
  Our use of the Fubini--Tonelli theorem in the sixth line is justified by Proposition \ref{pro:l1}.
\end{proof}

\end{document}
